% !TEX root = ../agr-paper.tex
% ---------------------------------------------------------------------
% Every limitation a reviewer could raise appears here first, in our
% words, at full strength. None is softened on the way across.
%
% THE LIST IS THE CONTRACT. scripts/check_paper_numbers.py enforces the
% roster by phrase, so nothing drops silently. The full set, and where
% each comes from:
%
%   1  sample size ......................... 400/dataset, +-4-5 pts
%   2  one backbone ....................... gpt-5.4-mini, frozen
%   3  nondeterminism ..................... ~67% trajectory stability;
%                                           the repeat matrix was
%                                           affordable and not run
%   4  environment friendlier ............. 97%+ reachability ceiling
%   5  ablation power ..................... n=200/198, three nulls,
%                                           one run, one at a time
%   6  scope .............................. English factoid, static,
%                                           no literals or ordinals
%   7  WRONGFUL ACCEPTANCE UNMEASURED ..... most serious item; also
%                                           kills auditability of the
%                                           output contract, because
%                                           the log keeps the COUNT of
%                                           supporting triples, not the
%                                           triples
%   8  view of the graph not held constant  candidate widths, Sec 5.3
%   9  static baseline is radius-bounded .. one logical hop, Sec 4.2
%  10  topic entities given, not linked ... setup gives only the benefit
%  11  measured accuracy has a floor ...... extraction bug on one shape
%  12  single-annotator adjudication ...... author, consensus visible
%  13  homonyms merge ..................... graph is name-keyed
%  14  judge missed its own bar ........... kappa 0.6995 vs 0.70
%  15  one relabelling after seeing result
%  16  structural grounding is not answer adequacy (the echo attractor)
%
% Added in the September 2026 reframe, without a roster phrase because
% each sits inside an existing item: the agentic comparator is one
% system and Plan-on-Graph was not run (8); the census reached 256 of
% 262 (12); the budget cap is a choice and the split conditions on an
% outcome (1/8).
%
% Close on the experiments the reframe makes necessary, in the order
% they would be run.
% ---------------------------------------------------------------------

\section{Discussion and Limitations}\label{sec:discussion}

The limits of this study bound its claims, and we state all of them
rather than the flattering subset. The first group bounds the
instrument, the second bounds the comparisons, and the third bounds the
measurements.

\subsection{What Bounds the Instrument}\label{sec:limits-contribution}

\textbf{Wrongful acceptance is unmeasured, and it is the most serious
gap in this work.} The verifier logs the claims it rejects and not the
ones it accepts. So a claim certified on an edge unrelated to the
relation it asserts leaves no trace. The exposure is real rather than
hypothetical, because both structural routes ignore the relation and the
direction. Any edge between two entities certifies any claimed relation
between them, so a claim that X is Y's mother is accepted on an edge
recording that X is Y's child. We therefore report one polarity at
census scale and the other anecdotally, and that asymmetry is a property
of the logging decision rather than of the verifier. There is no way,
from the committed record, to know how many wrongful acceptances sit
inside the $761$ questions that ended grounded.

The exposure is best stated in the claim unit, where its scale is
visible. On test the verifier decomposed $2{,}110$ claims and accepted
$2{,}008$ of them. The graph-wide adjacency probe accounts for $39$ of
those acceptances. The rest were made either by traversed adjacency,
which is equally relation-blind, or by the entailment check, which is
not, and the record does not say which. So the share of accepted claims
certified without any test of the asserted relation lies somewhere in
$[39, 2{,}008]$. That interval is two orders of magnitude wide, and we
report it as one rather than choose a point inside it. The same gap
means the layer has no measured precision or recall as a classifier of
claims, which is the first number a verification component should carry.

\textbf{The output contract cannot be audited from the run records.} AGR
emits an answer paired with the traversed triples supporting it, and the
pairing does hold inside the run: the answerer receives the verifier's
triple list and emits against it. But the logger writes the \emph{count}
of supporting triples and discards the list, so no committed artifact in
this work contains a single supporting triple. On test, $12$ of AGR's
$367$ answered WebQSP questions and $9$ of its $308$ on
ComplexWebQuestions were emitted with a count of zero, which is the
contract's own measure of how often it delivers nothing. A reader can
confirm that the other answers came from a system that tracked its
evidence. They cannot inspect that evidence. Persisting accepted claims
with their route and their matching triples is one logging change, and
it would repair this limitation and the previous one together.

\textbf{Structural grounding is not answer adequacy.} Both navigating
systems assert zero ungrounded entities, and under the definition used
that is close to guaranteed for any system that copies entities out of
tool results. It says their answers are made of the graph's own
material. It does not say they answer the question. The echo attractor
of \Cref{sec:echo} is the counter-example: a wrong answer that passes
every check this paper applies. The verification layer offers no
protection against it.

\subsection{What Bounds the Comparisons}\label{sec:limits-comparison}

\textbf{Sample size.} Every result rests on $400$ questions per dataset
rather than the full splits, giving $95\%$ bootstrap intervals of
roughly $\pm4$ to $\pm5$ points on Hits@1. The main accuracy margins
clear that width: $9.5$ points over the agentic baseline on WebQSP and
$8.7$ on ComplexWebQuestions. So does the planner effect, which is a
paired test on per-question outcomes rather than a comparison of two
intervals. The groundedness contrast, at zero against $22.1\%$, is not
close. What the sample does not support is fine distinctions. The
semantic-tier gaps of \Cref{sec:groundedness}, the hedge-rate difference
on ComplexWebQuestions, and the last step of the rising hop-depth curve
are within noise, and we report them as such.

\textbf{The margin is a statement about one regime.} The $25$-call cap
was pre-specified, but it is a choice, and \Cref{sec:margin} shows that
the entire margin over the agentic baseline lives on the questions where
that cap truncates it or where the baseline's candidate cut discarded a
gold-path relation. Where neither binds, the baseline is ahead. The
split that shows this conditions on the baseline's own outcome, and
within hop strata the clipped half is as hard as the finished half on
WebQSP and harder on ComplexWebQuestions. A cap sweep would move the
boundary, and the paper reports the clip rates so a reader can see how
far, but it did not run the sweep.

\textbf{An identical environment is not identical access to it.} The
agentic baseline truncates candidate sets at $40$ relations and $20$
neighbours against AGR's $300$ and $200$. That cut binds on about a
third of its expansions and effectively none of ours, and on $89$ and
$51$ questions it discarded the relation continuing a shortest path to a
gold answer (\Cref{sec:margin}). The comparison holds the graph constant
and does not hold the view of it constant, which is why the baseline's
figures on the questions it finishes are a lower bound rather than an
estimate, and why the aggregate margin is not read as an architectural
one. The agentic comparator is also one system. Plan-on-Graph, whose two
mechanisms AGR's planner and backtracker re-host
(\Cref{sec:rel-agentic}), was not run, so the planner result speaks to
decomposition as implemented here and to Plan-on-Graph's by analogy
only.

\textbf{The static graph baseline is radius-bounded.} Its
one-logical-hop expansion cannot contain a two-hop chain, so its decay
across hop strata measures the radius as much as the paradigm, and
\Cref{sec:hops} declines to offer it as evidence for either. Its
unordered $100$-edge fanout cap compounds this.

\textbf{One backbone.} All five systems run on \textsf{gpt-5.4-mini},
frozen and closed. This paper establishes how these architectures
compare \emph{on one backbone}, not how the comparison scales across
model capacities.

\textbf{Nondeterminism.} Temperature-zero decoding against a hosted API
is greedy, not deterministic, and measured trajectory stability is about
$67\%$. Each system and each ablation condition was run once. No seeds
were averaged, and a reader should not assume they were. The intervals
and paired tests resample questions, so they bound the sampling variance
of the question draw and nothing else; each system's answer set enters
them as fixed. Two things make a rank reversal from decoding noise alone
unlikely: the defended margins are large relative to their intervals,
and the paired tests are computed per question. But unlikely is the
strongest statement available. The repeated matrix that would bound the
variance was affordable at the observed cost and was not run, so this is
a gap in the design rather than a budget constraint.

\textbf{Ablation power.} The four ablations run on half-splits of $n =
200$ and $n = 198$, and the paired test sees only the questions on which
two conditions disagree, as few as one for claim verification. Three of
four components show \emph{no effect detected at this power}, which is
not the same as no effect. \Cref{sec:power} gives the smallest
difference each condition could have detected, between $4.0$ and $5.5$
accuracy points where a test was possible at all. The components were
also removed one at a time, so the three nulls do not say whether a
system with none of them would match this one; that condition was not
run.

\subsection{What Bounds the Measurements}\label{sec:limits-measurement}

\textbf{The environment is friendlier than the knowledge base it stands
in for.} Per-question subgraphs with a $97\%$-plus answer-reachability
ceiling are a more forgiving substrate than full Freebase.
Between-system comparisons hold, because all five systems navigate the
identical graph. Absolute accuracies do not transfer, and \Cref{sec:rog}
shows the size of the gap to a published system that was also trained on
the benchmarks.

\textbf{Topic entities are given, not linked.} Both graph-navigating
systems seed from the datasets' annotated topic entities. That isolates
navigation from entity linking, which is the intent. But it means every
reported accuracy assumes a linking step a deployed system would have to
perform, and which would introduce errors of its own.

\textbf{Entities are keyed by name, so homonyms merge.} Distinct
real-world entities sharing a surface form collapse to one node. That
can manufacture a path between things not actually connected, letting a
wrong answer verify as grounded, and it can equally destroy a
distinction a question depends on. Neither direction is measured.
Structural groundedness is therefore a claim about this graph's node
identity rather than about the world.

\textbf{The measured accuracy has an unmeasured floor.} An extraction
defect converts correct reasoning into a wrong entity list on one
specific question shape. On ``list the professions'' and ``what did $X$
do'' questions the drafted text names every gold value and the scored
entity list collapses to the sentence's subject. Accuracy on that shape
is a lower bound rather than an estimate.

\textbf{The judge missed its own bar, and the adjudication was
single-annotator.} Agreement between the semantic-tier judge and a human
annotator reached $\kappa = 0.6995$ against a pre-specified $0.70$. That
is a miss, by $0.0005$. It is why those numbers are reported as
corroborating rather than decisive, and why no conclusion rests on them
alone. The judge also runs on the same backbone as the systems it
judges. Separately, the benchmark-defect exclusions and the failure
census rest on single-annotator judgements made by an author with the
consensus signal visible, and the census reached $256$ of its $262$
questions (\Cref{sec:erroranalysis}). One development-set relabelling
was also made after its outcome was known. It is flagged in place, and
no aggregate was rescored.

\textbf{Scope.} English factoid questions over a static graph, with no
temporal facts, no graph construction, and an offline prototype. The
environment carries no typed literals and the tool API no ordering
operator, so questions answered by a date, a quantity, or a superlative
are out of reach by construction (\Cref{sec:environment}). Wall-clock
figures are reported as magnitudes with their cache caveat
(\Cref{sec:cost}), not as a production-readiness claim.

\subsection{Outlook}\label{sec:outlook}

The findings above fix the order of the next experiments, and each is
cheap against the cost of the study it would complete. First, the
agentic baseline at AGR's candidate widths under a sweep of the call
cap, with the budget split of \Cref{tab:togsplit} reported at each
setting, which tests whether the affordability reading survives equal
access to the graph. Second, replication: every system and every
ablation condition run several times, so that a delta of a few
hundredths of F1 can be read against the run-to-run variance rather than
assumed to exceed it, together with the joint-removal condition that
\Cref{sec:power} names. Third, the full test splits, which quadruple the
discordant pairs the paired tests see. Fourth, one logging change:
persist each accepted claim with the route that accepted it and the
triples that matched. That converts the two limitations opening this
section into measurable quantities, and it makes the verifier's
precision and recall a number a labelled sample can produce.

The most actionable finding points forward rather than back. If
decomposition helps only when a question has compositional structure,
improving accuracy by $0.083$ F1 where that structure is absent and
costing $0.047$ where it is present, then the planner should not run
unconditionally. A hop-count predictor gating the planner would capture
both arms of that asymmetry, and it can be evaluated with the ablation
instrument used here. More generally, the result suggests that agentic
components should be treated as conditionally rather than universally
beneficial, and that the condition is worth measuring.
